\section{Methods}
\label{sec:methods}

This section describes the empirical testbed (Section~\ref{subsec:testbed}), the metrics that summarise each run (Section~\ref{subsec:metrics}), and the experimental blocks (Section~\ref{subsec:overview}) that answer the research questions of Section~\ref{sec:proposed_approach}. Each block maps to one research question and to a prediction stated earlier.

\subsection{Experimental Setup}
\label{subsec:setup}

\subsubsection{Testbed}
\label{subsec:testbed}

The experiments run on two testbeds of different scale. A small, fast testbed, rotated MNIST on a two-layer MLP, carries the dense mechanistic sweeps: the driver block, the curriculum sweep, the asymmetric-PER damping sweep, and the momentum cross. A larger testbed, a domain-shift CIFAR-10 on a CIFAR-adapted ResNet-18, tests whether the findings survive a deeper backbone and a stronger task shift, the regime where the local condition behind the bound of Section~\ref{sec:theory} weakens. Both testbeds share one design choice: the label space and task structure are held fixed and only the input distribution shifts, and the headline sequences use two tasks, so there is exactly one transition and exactly one instance of the stability gap.

\paragraph{Datasets.}

The small benchmark is rotated MNIST (rot-MNIST). Following \citet{delange2023continualevaluationlifelonglearning}, the first task $T_0$ presents the digits at $-90^\circ$ and the second task $T_1$ applies a further $90^\circ$ rotation. Each task uses the standard $60{,}000$ training and $10{,}000$ test images. The $90^\circ$ rotation separates the two input distributions well while keeping the tasks structurally identical.

The large benchmark is a domain-shift CIFAR-10: $T_0$ uses clean CIFAR-10 and $T_1$ applies the Gaussian-noise corruption at severity $3$ \citep{hendrycks2019robustness}. The ten classes are shared, so only the input distribution shifts. A three-task variant ($[\text{none}, \text{Gaussian noise}, \text{shot noise}]$) and two alternative corruptions are used in the generalisation block of Section~\ref{subsec:cifar}.

\paragraph{Models.}

The rot-MNIST sweeps use a two-layer multilayer perceptron (MLP) with ReLU activations (Table~\ref{tab:mlp}); the CIFAR-10 block uses a CIFAR-adapted ResNet-18 (Table~\ref{tab:resnet}). The MLP is cheap to train, which matters because the sweeps span many conditions and seeds, and it is the setting where the local condition of Section~\ref{sec:theory} is most plausible. The ResNet-18 is the deeper, more realistic counterpart.

\begin{table}[htbp]
    \centering
    \small
    \begin{tabular}{@{}l p{0.55\columnwidth}@{}}
        \hline
        Component        & Value                                    \\
        \hline
        Input            & $784$ (flattened $28\times28$ greyscale) \\
        Hidden layers    & $2 \times 400$ (fully connected, ReLU)   \\
        Output           & $10$-class logits                        \\
        Regularisation   & none                                     \\
        Total parameters & $478{,}410$                              \\
        \hline
    \end{tabular}
    \caption{MLP backbone used for the rot-MNIST sweeps.}
    \label{tab:mlp}
\end{table}

\begin{table}[htbp]
    \centering
    \small
    \begin{tabular}{@{}l p{0.55\columnwidth}@{}}
        \hline
        Component        & Value                                                                                               \\
        \hline
        Input            & $3\times32\times32$ RGB                                                                             \\
        Stem             & $3\times3$ conv, stride $1$, no max-pool (CIFAR adaptation)                                         \\
        Stages           & $4$ residual stages, BasicBlocks $[2,2,2,2]$, widths $[64,128,256,512]$                             \\
        Normalisation    & BatchNorm (default); GroupNorm ($\min(32, C)$ groups) variant crossed in Section~\ref{subsec:cifar} \\
        Head             & global average pool $+$ linear, $10$-class logits                                                   \\
        Total parameters & ${\approx}11.2$M                                                                                    \\
        \hline
    \end{tabular}
    \caption{CIFAR-adapted ResNet-18 backbone used for the CIFAR-10 block. The $3\times3$ stride-$1$ stem and removed max-pool keep the feature map at $32\times32$; all other details follow the ResNet-18.}
    \label{tab:resnet}
\end{table}

\paragraph{Training protocol.}

Table~\ref{tab:training} lists the training hyperparameters. rot-MNIST runs in the online regime of one epoch per task, standard for stability-gap analysis: it gives a single clean pass through the transition. The ResNet-18 needs several passes to converge before the transition, so the CIFAR block uses ten epochs per task. Everything else (optimiser, learning rate, batch size) is held common across the two benchmarks.

\begin{table}[htbp]
    \centering
    \small
    \begin{tabular}{@{}l p{0.27\columnwidth} p{0.27\columnwidth}@{}}
        \hline
        Hyperparameter  & rot-MNIST (MLP)         & CIFAR-10 (ResNet-18)          \\
        \hline
        Epochs per task & $1$ (online)            & $10$                          \\
        Batch size      & $256$                   & $256$                         \\
        Optimiser       & SGD                     & SGD                           \\
        Learning rate   & $0.1$                   & $0.1$                         \\
        Momentum        & $0.0$ or $0.9$          & $0.9$ ($0.0$ in cross)        \\
        Weight decay    & $0.0$                   & $0.0$                         \\
        Replay buffer   & $1{,}000$ or $60{,}000$ & $5{,}000$                     \\
        Seeds           & $5$                     & $5$ or $3$                    \\
        \hline
    \end{tabular}
    \caption{Training protocol for both benchmarks.}
    \label{tab:training}
\end{table}

\subsubsection{Evaluation Protocol}
\label{subsubsec:eval_protocol}

Throughout training, every task's test accuracy is evaluated every $10$ steps. At each task transition the cadence increases to every step, so the dip is recorded at full resolution where it occurs. On rot-MNIST the per-step cadence covers the whole new task (one online epoch, ${\approx}234$ steps). On CIFAR-10 it is held for a $250$-step window after the transition and then drops back to every $10$ steps. A stability-gap tracker snapshots the pre-task accuracy on each past task before the transition and records the resulting (step, accuracy) series, from which it computes the scalar metrics below. A gradient tracker records, at the same cadence, the fidelity of the replay-buffer gradient against the exact past-task gradient; it is enabled for the driver block only, because each diagnostic step costs a full pass over the past-task training set.

\subsubsection{Metrics}
\label{subsec:metrics}

All metrics are recorded automatically and aggregated later. Unless noted otherwise, each is reported as the seed mean $\pm$ sample standard deviation; \texttt{gap\_area\_end} is reported as a seed mean only. Complete definitions of every logged metric are collected in Appendix~\ref{app:metric_defs}.

\paragraph{Stability-gap metrics.}
\label{subsubsec:gap_metrics}

The definition of the gap in Section~\ref{subsec:reference}, the transient excess below the level the task settles at, translates directly into the two primary metrics. Let $a_j^{\text{pre}}$ denote the baseline accuracy on past task $j$ before the switch and $a_j(t)$ its accuracy at step $t$ of the new task. The first metric, \texttt{gap\_depth}, is the maximum instantaneous excursion,
\begin{equation*}
    \texttt{gap\_depth} = \max_j \left( a_j^{\text{pre}} - \min_t a_j(t) \right),
\end{equation*}
which Section~\ref{subsec:overshoot} predicts to be overshoot-dominated. The second, \texttt{gap\_area\_end}, is the time-integral of the excursion below the settled level: with the reference point $b_j = \min\!\left(a_j^{\text{pre}},\, \bar{a}_j\right)$, where $\bar{a}_j$ is the stable accuracy the task eventually reaches, it integrates the shortfall $\max\!\left(0,\, b_j - a_j(t)\right)$ over the new task and sums across past tasks. Anchoring to the recovered level implements the split of Section~\ref{subsec:reference}: a task that steps down to a permanently lower plateau contributes its loss to the standard forgetting metrics and approximately zero to the area, while a true dip-and-recover transient registers fully. The area is where the arc of Section~\ref{subsec:trajectory_problem} and the lingering stochastic residual appear. A third metric, \texttt{recovery\_steps}, counts the steps until every past task climbs back to within $90\%$ of its pre-task baseline, and stays undefined when accuracy never drops that far.

\paragraph{Continual-learning metrics.}

To place the gap in the standard continual-learning frame, we report the continual-evaluation metrics of \citet{delange2023continualevaluationlifelonglearning}, computed from the task-boundary accuracy matrix logged for every run. The two that carry the argument are average final accuracy (ACC), which checks that a gate buys a lower gap at preserved accuracy, and the worst-case retention measures min-ACC and WC-ACC. FORG and the windowed measures $\mathrm{WF}_w$/$\mathrm{WP}_w$ are reported for completeness; $\mathrm{WP}_{100}$ serves as the short-horizon plasticity readout in the gate comparisons. Definitions are in Appendix~\ref{app:metric_defs}.

\paragraph{Gradient-fidelity diagnostics.}
\label{subsubsec:grad_metrics}

For the driver block, the gradient tracker compares the replay-buffer gradient $g_{\text{replay}}$ against the exact empirical past-task gradient $g_{\text{true}}$ (a separate full-data pass). Two diagnostics quantify how faithfully a finite buffer approximates the past data \citep{aljundi2019gradientbasedsampleselection}: the directional fidelity $\cos(g_{\text{replay}}, g_{\text{true}})$ and the magnitude fidelity $\|g_{\text{replay}}\|/\|g_{\text{true}}\|$. These supply the per-step evidence for the variance driver of Section~\ref{subsec:overshoot}. A third ratio, $\|g_{\text{replay}}\|/\|g_{\text{new}}\|$, measures the boundary imbalance; it is also the feedback signal of the adaptive curriculum (Section~\ref{subsec:curriculum}).

\paragraph{Statistical comparison.}

Across conditions we compare metrics with the paired Wilcoxon signed-rank test on the seed-paired difference. Seeds are paired by index, so the per-seed deltas factor out the seed-induced variance that would otherwise dominate a five-seed sample. We prefer Wilcoxon over the paired $t$-test because \texttt{gap\_depth} and accuracy are bounded in $[0,1]$ and typically right-skewed near zero, so the $t$-test's normality assumption is suspect. With five paired seeds the smallest possible two-sided $p$-value is $0.0625$, reached exactly when all five differences share a sign; we therefore read $p = 0.0625$ as complete directional agreement and let effect sizes separate large results from marginal ones. For equivalence-style claims we report the seed-paired difference with its bootstrap $95\%$ confidence interval.

\subsection{Experiments}
\label{subsec:overview}

\begin{table*}[htbp]
    \centering
    \small
    \begin{tabular}{@{}p{3.0cm} p{1.0cm} p{0.9cm} p{5.9cm} p{3.4cm}@{}}
        \hline
        Block                        & Section                 & RQ        & Prediction under test                                                                                          & Conditions                                                                          \\
        \hline
        Overshoot drivers            & \ref{subsec:drivers}    & RQ1       & Removing the drivers flattens the drop; the arc survives in the area                                           & G1--G4                                                                              \\
        $\lambda$-curriculum sweep   & \ref{subsec:curriculum} & RQ2a--c   & Depth falls with $N$; adaptive matches without the knob; $\lambda_{\min}$ trades depth for plasticity          & C1--C7                                                                              \\
        Full-buffer diagnostic       & \ref{subsec:curriculum} & RQ1, RQ2d & Curriculum $+$ exact gradients close the gap without momentum                                                  & C8, C8$_M$                                                                          \\
        Momentum cross               & \ref{subsec:momentum}   & RQ2d, RQ3 & Momentum composes with source attenuation and undoes per-step attenuation                                      & every rot-MNIST condition at $\mu \in \{0.0, 0.9\}$                                 \\
        Asymmetric-PER $\delta$ sweep & \ref{subsec:per_asym}  & RQ3       & Depth and arc shrink as $\delta \to 0$ at little plasticity cost; feedforward failure modes appear             & $\delta \in \{1.0, 0.3, 0.1, 0.03\}$ $\times$ buffer legs, plus a solver control    \\
        CIFAR-10 generalisation      & \ref{subsec:cifar}      & RQ4       & Both gates and the momentum predictions carry to a deeper backbone                                             & vanilla, curriculum, asym.\ PER $\times$ \{BN, GN\}; momentum cross; novel corruptions \\
        \hline
    \end{tabular}
    \caption{Overview of the experimental blocks, the research question each answers, and the prediction each tests.}
    \label{tab:overview}
\end{table*}

Each block tests one or more predictions of Sections~\ref{sec:motivation} and~\ref{sec:theory} (Table~\ref{tab:overview}). All conditions use five seeds and the protocol of Section~\ref{subsec:testbed} unless noted otherwise.

\subsubsection{Overshoot Drivers}
\label{subsec:drivers}

\paragraph{Aim.} Remove the overshoot drivers of Section~\ref{subsec:overshoot} by construction, one at a time, and observe what remains. This supplies the driver half of RQ1; the path half comes from the curriculum conditions, and the two halves meet in the analysis.

\begin{table*}[htbp]
    \centering
    \small
    \begin{tabular}{p{2.9cm} p{2.4cm} p{2.7cm} p{6.2cm}}
        \hline
        Condition                   & Replay buffer    & Gradient handling  & Driver removed by construction                          \\
        \hline
        G1 --- Vanilla ER           & $1$k reservoir   & unbalanced         & none (all drivers active)                               \\
        G2 --- Full-data ER         & $60$k (all past) & unbalanced         & estimator variance (exact replay gradient)              \\
        G3 --- Balanced ER          & $1$k reservoir   & unit-norm balanced & boundary imbalance                                      \\
        G4 --- Full-data + balanced & $60$k            & unit-norm balanced & imbalance and variance                                  \\
        \hline
    \end{tabular}
    \caption{Driver-removal conditions. Each condition switches off one or two overshoot drivers relative to vanilla ER; what all four leave in place is the descent path itself, so the arc should survive throughout.}
    \label{tab:drivers}
\end{table*}

\paragraph{Design.} Table~\ref{tab:drivers} lists the four conditions. The balanced conditions G3 and G4 compute the replay and new-task gradients in two separate backward passes, normalise each to unit norm, and combine them $50/50$, so the two terms enter the update with equal magnitude from the first step. This switches off the boundary imbalance: at $\theta_0^*$ an unbalanced step is driven almost entirely by $g_{\text{new}}$, while a balanced step carries the replay direction at equal weight. The full-data conditions G2 and G4 hold the entire past-task training set in the buffer and compute the replay gradient on every stored sample at every step. The replay gradient becomes the exact empirical past-task gradient, which switches off the estimator variance; vanilla ER's $256$-sample draw is a noisy estimator of the same quantity.

\paragraph{How the conditions are read.} The drivers interact (Section~\ref{subsec:overshoot}), so we read the conditions as interventions rather than as an additive decomposition. Depth differences between conditions measure how much of the drop each driver carries in that context. The area, and the shape of the per-step curve, carry the level-one signal: with variance switched off (G2, G4) any dip-and-recover bend that remains in the curve is deterministic, which is the arc's fingerprint.

\paragraph{Prediction.} Depth falls as drivers are removed and reaches its floor at G4. The arc survives every removal: G2 and G4 show a smooth, deterministic bend in the per-step curve and a clear residual area, because balancing and exact gradients leave the descent path bent. Only a path-level intervention should remove that bend, which the curriculum conditions test.

\subsubsection{The \texorpdfstring{$\lambda$}{lambda}-Curriculum Sweep}
\label{subsec:curriculum}

\paragraph{Aim.} Test the scalar feedback gate: the $O(1/N)$ prediction of Equation~\eqref{eq:bound} (RQ2a), the adaptive schedule (RQ2b), the $\lambda_{\min}$ floor (RQ2c), and, with the full-buffer diagnostic, the claim that the curriculum removes the arc and leaves only stochastic overshoot (RQ1).

\paragraph{Design.} All curriculum conditions run on top of standard ER ($1$k reservoir buffer, unbalanced gradients), so the schedule is the only changed variable (Table~\ref{tab:curriculum}).

\begin{table*}[htbp]
    \centering
    \small
    \begin{tabular}{@{}p{0.9cm} p{3.4cm} p{4.8cm} p{3.4cm} p{1.4cm}@{}}
        \hline
        Label & Condition                           & $\lambda(t)$ schedule                                                                  & Hyperparameter                           & Tests     \\
        \hline
        ---   & Vanilla ER (baseline)               & $\lambda \equiv 1$ from step $1$ (abrupt switch)                                       & ---                                      & reference \\
        C1    & Linear $N{=}50$                     & $\mathrm{clip}(t/50,\, 0,\, 1)$                                                        & $N = 50$                                 & RQ2a      \\
        C2    & Linear $N{=}100$                    & $\mathrm{clip}(t/100,\, 0,\, 1)$                                                       & $N = 100$                                & RQ2a      \\
        C3    & Linear $N{=}200$                    & $\mathrm{clip}(t/200,\, 0,\, 1)$                                                       & $N = 200$                                & RQ2a      \\
        C4    & Adaptive                            & $\mathrm{clip}\!\left(\mathrm{EMA}_\alpha(r),\, 0,\, 1\right)$                         & $\alpha = 0.05$                          & RQ2b      \\
        C5    & Adaptive $+\,\lambda_{\min}{=}0.05$ & $\max\!\left(\lambda_{\min},\, \mathrm{clip}(\mathrm{EMA}_\alpha(r),\, 0,\, 1)\right)$ & $\alpha = 0.05$, $\lambda_{\min}{=}0.05$ & RQ2c      \\
        C6    & Adaptive $+\,\lambda_{\min}{=}0.10$ & as C5                                                                                  & $\lambda_{\min}{=}0.10$                  & RQ2c      \\
        C7    & Adaptive $+\,\lambda_{\min}{=}0.20$ & as C5                                                                                  & $\lambda_{\min}{=}0.20$                  & RQ2c      \\
        C8    & C7 $+$ full $60$k buffer            & as C7, exact replay gradient                                                           & diagnostic                               & RQ1, RQ2d \\
        \hline
    \end{tabular}
    \caption{$\lambda$-curriculum conditions, all applied on top of standard ER. $\mathrm{EMA}_\alpha$ is an exponential moving average with smoothing $\alpha$, and $r := \|g_{\text{replay}}\| / \|g_{\text{new}}\|$ is the gradient-norm ratio. Each condition is run at momentum $0.0$ and $0.9$.}
    \label{tab:curriculum}
\end{table*}

The adaptive schedule (C4) is the feedback variant: it sets $\lambda$ to a clipped exponential moving average of the gradient-norm ratio $r := \|g_{\text{replay}}\| / \|g_{\text{new}}\|$. At $\theta_0^*$ the replay gradient is near zero by stationarity, so $r$ and $\lambda$ start near $0$. As the iterate moves, $\|g_{\text{replay}}\|$ grows while $\|g_{\text{new}}\|$ shrinks, so $\lambda \to 1$ at the pace the damage signal dictates, with no ramp length to choose. The floor $\lambda_{\min}$ (C5--C7) keeps the schedule from sitting at $\lambda \approx 0$ during the first steps, when the EMA of $r$ is still close to zero, and thereby buys back early plasticity.

\paragraph{Full-buffer diagnostic (C8).} C8 combines the shipping curriculum (C7) with the full $60$k buffer of G2/G4, so the path is fixed \emph{and} the variance driver is switched off at the source. Its role is diagnostic: if the curriculum's remaining transient at momentum $0$ is stochastic overshoot, C8 should close the gap with no momentum at all, and the comparison of C8 with C8$_M$ separates momentum's two candidate roles, variance reduction (replaceable by the full buffer) and acceleration (which only momentum provides). C8 stores the entire past task and therefore breaks the limited-memory premise of continual learning; it is a probe, and C7 remains the practical configuration.

\paragraph{Prediction.} Depth falls as $N$ grows (RQ2a) while ACC stays near vanilla ER; the adaptive schedule reaches at least the best linear depth without the knob (RQ2b); raising $\lambda_{\min}$ buys early $T_1$ speed and final accuracy for a modest depth increase (RQ2c). For C8, the account predicts the lowest gap of any condition at momentum $0$, together with reduced plasticity, because all three suppressors of the new task (the schedule, the strong exact anchor, and the missing accelerator) then stack.

\subsubsection{Momentum Cross}
\label{subsec:momentum}

Every rot-MNIST condition, the driver G-series, the curriculum C-series, and the asymmetric-PER sweep, is run at momentum $0.0$ and $0.9$; momentum-on runs carry an $M$ subscript (e.g.\ $\text{G1}_M$). The cross tests the composability prediction of Section~\ref{subsec:gates}: whether momentum helps or hurts a method's gap should depend on where the method attenuates the push.

Three cases follow from the account. For vanilla ER, the drop is set by the deterministic, unopposed first step; that step carries no accumulated velocity, so momentum should leave the \emph{depth} unchanged while it smooths the recovery and shrinks the area. For the curriculum, the push is attenuated at source: momentum finds a small, clean signal, so it should average away the stochastic residual (closing depth and area together) and restore final accuracy, the composition of RQ2d. For asymmetric PER, the push is attenuated per step while the underlying gradient keeps its size: accumulation across steps can rebuild what the filter removes, so momentum should partially undo the protection and re-inflate the depth, at a simultaneous accuracy benefit. Accuracy itself should improve under momentum for every method, since acceleration is the plasticity lever regardless of gate.

\subsubsection{Asymmetric-PER \texorpdfstring{$\delta$}{delta} Sweep}
\label{subsec:per_asym}

\paragraph{Aim.} Test the directional feedforward gate (RQ3): does the gap shrink as the filter strengthens, at what plasticity cost, and do the two predicted feedforward failure modes appear?

\paragraph{Method.} Asymmetric PER implements Equation~\eqref{eq:asym_per}. The metric $F_{\text{rep}}$ is the Fisher information of the replay batch, the model's own estimate of which directions the past task depends on \citep{amari1998natural}. The filtered push $\delta\,(F_{\text{rep}}+\delta I)^{-1} g_{\text{cur}}$ is computed each step by warm-started conjugate gradients over Fisher-vector products ($25$ iterations), so no matrix is ever formed. The normalisation by $\delta$ keeps flat directions at unit gain: the gain along an eigendirection with curvature $\rho$ is $\delta/(\rho+\delta)$, so directions the past task is flat in pass through whole, and directions it is steep in are shrunk toward zero. The damping $\delta$ is the single knob. At $\delta \to \infty$ every direction passes and plain ER returns; as $\delta \to 0$ the push concentrates in the replay-flat directions while the restoring gradient $g_{\text{rep}}$ keeps full strength throughout. Implementation details are in Appendix~\ref{app:implementation}.

\paragraph{Design.} The sweep crosses the damping $\delta \in \{1.0, 0.3, 0.1, 0.03\}$ with two buffer legs, at both momentum settings (Section~\ref{subsec:momentum}). The \emph{standard} leg uses the $1$k reservoir buffer and matches the practical regime of the curriculum sweep. The \emph{full-buffer} leg computes the replay loss gradient over all $60$k past samples, which switches off the variance driver exactly as in G2/G4; the Fisher stays estimated on a sampled batch, since a full-buffer Fisher-vector product per CG iteration would dominate the runtime. The full-buffer leg isolates the mechanism: with the stochastic overshoot gone, the remaining area measures the deterministic arc, so the sweep reads out directly whether the directional gate shrinks the arc itself. A solver control reruns the strongest filter ($\delta = 0.03$, full buffer) at $60$ CG iterations; it separates properties of the gate from artefacts of an under-converged solve.

\paragraph{Prediction.} Depth and area fall monotonically as $\delta \to 0$, on both buffer legs, at little cost in ACC and $\mathrm{WP}_{100}$: the filter blocks the push only where the past task is steep, so plasticity should survive. The taxonomy of Section~\ref{subsec:gates} predicts where this breaks. The sampled Fisher can rate a direction as flat that the true past-task loss punishes further out; at strong filtering, new-task pressure should escape through such blind spots and produce localised failures that a feedback signal would have caught. And under momentum the per-step protection should partially unravel (Section~\ref{subsec:momentum}).

\subsubsection{Generalisation: CIFAR-10}
\label{subsec:cifar}

\paragraph{Aim.} Test whether both gates and the momentum predictions carry to a deeper backbone and a stronger task shift (RQ4).

\paragraph{Protocol.} All CIFAR conditions use the ResNet-18 of Table~\ref{tab:resnet}, ten epochs per task, batch size $256$, learning rate $0.1$, and a replay buffer of $5{,}000$ samples, with the evaluation scheme of Section~\ref{subsubsec:eval_protocol}. The headline conditions run at momentum $0.9$, the standard setting for this backbone. Each method is crossed with BatchNorm and GroupNorm, because the two normalisations respond differently at the switch: the input shift disrupts BatchNorm's running statistics, which must re-adapt during the first steps of the new task, while GroupNorm uses no cross-batch statistics, so its transient reflects optimisation dynamics alone.

\begin{table}[H]
    \centering
    \small
    \begin{tabular}{@{}p{0.7cm} p{3.2cm} p{1.5cm} p{0.8cm}@{}}
        \hline
        Label & Method                            & Norm      & $\mu$  \\
        \hline
        D1    & Vanilla ER                        & BatchNorm & $0.9$  \\
        D3    & Adaptive $\lambda_{\min}{=}0.20$  & BatchNorm & $0.9$  \\
        D4    & Vanilla ER                        & GroupNorm & $0.9$  \\
        D6    & Adaptive $\lambda_{\min}{=}0.20$  & GroupNorm & $0.9$  \\
        P1    & Asym.\ PER $\delta{=}0.1$         & BatchNorm & $0.9$  \\
        P2    & Asym.\ PER $\delta{=}0.1$         & GroupNorm & $0.9$  \\
        \hline
        M1    & Vanilla ER                        & GroupNorm & $0.0$  \\
        M2    & Adaptive $\lambda_{\min}{=}0.20$  & GroupNorm & $0.0$  \\
        M3    & Asym.\ PER $\delta{=}0.1$         & GroupNorm & $0.0$  \\
        \hline
    \end{tabular}
    \caption{CIFAR-10 conditions (two-task Gaussian-noise shift, buffer $5{,}000$, five seeds). Top: the headline block at momentum $0.9$. Bottom: the momentum cross, which pairs momentum-$0$ runs of all three methods against their momentum-$0.9$ counterparts at GroupNorm.}
    \label{tab:cifar}
\end{table}

\paragraph{Headline block.} Table~\ref{tab:cifar} (top) compares the three methods at both normalisations. The curriculum ships as the adaptive schedule at $\lambda_{\min} = 0.20$, the configuration selected on rot-MNIST; asymmetric PER ships at $\delta = 0.1$, a strong rot-MNIST setting on both axes that stays clear of the blind-spot regime of the sweep. Both methods carry their rot-MNIST configurations over without a CIFAR-specific sweep, so they meet the deeper backbone on equal footing.

\paragraph{Momentum cross.} Table~\ref{tab:cifar} (bottom) adds momentum-$0$ runs of all three methods at GroupNorm, whose transient is free of normalisation-statistics effects. Together with the momentum-$0.9$ rows, this repeats the three-way composability test of Section~\ref{subsec:momentum} on the deep backbone: depth unchanged for vanilla, composition for the curriculum, re-inflation for asymmetric PER. Momentum also drives convergence on this backbone, so accuracy is compared within a momentum setting rather than across.

\paragraph{Generalisation block.} A further block tests whether the curriculum's result depends on the particular corruption or on having a single transition. It pairs the shipping curriculum against a vanilla-ER control on a two-task shot-noise shift, a two-task contrast shift, and a three-task sequence $[\text{none}, \text{Gaussian noise}, \text{shot noise}]$, each at both normalisations over three seeds (Table~\ref{tab:cifar_gen}). The paired control keeps the comparison interpretable: the question is whether the gate still beats vanilla on data it was never tuned on.

\begin{table}[H]
    \centering
    \small
    \begin{tabular}{@{}p{3.0cm} p{1.0cm} p{2.6cm}@{}}
        \hline
        Shift                         & Tasks & Conditions        \\
        \hline
        Shot noise                    & $2$   & ship vs.\ vanilla \\
        Contrast                      & $2$   & ship vs.\ vanilla \\
        Gaussian noise $+$ shot noise & $3$   & ship vs.\ vanilla \\
        \hline
    \end{tabular}
    \caption{CIFAR-10 generalisation block. Each shift runs the shipping curriculum (adaptive $\lambda_{\min}{=}0.20$) and a paired vanilla-ER control, at both normalisations, three seeds.}
    \label{tab:cifar_gen}
\end{table}

\paragraph{Prediction.} Qualitative carry-over of the rot-MNIST results. The curriculum reduces \texttt{gap\_depth} and \texttt{gap\_area\_end} below vanilla ER at preserved ACC under both normalisations, and keeps beating its control on the novel corruptions. Asymmetric PER lands between vanilla and the curriculum in depth at momentum $0.9$, its weakest regime, and improves at momentum $0$. The momentum cross reproduces the three composability signatures. On this backbone the transient stays open for many steps, so the area carries as much of the comparison as the depth; the quantitative $O(1/N)$ scaling is a local result and is checked only qualitatively here.
